\setcounter{chapter}{5}
\chapter{Carathéodory流の測度論}

\paragraph{本章の目標}
Lebesgue可測性は，内測度と外測度の一致として定義した
（\cref{def:lebesgue-measurable-local}）．
一方で，同じ可測性はLebesgue外測度に対するCarathéodory条件でも特徴づけられる
（\cref{cor:lebesgue-caratheodory-equivalence}）．
本章ではこの後者の見方を一般化し，任意の外測度から可測集合と測度を取り出す方法，
および半加法族上の体積から測度を作る方法を整理する．

\section{動機づけ}

Lebesgue内測度は，コンパクト集合で内側から近似する量として導入した．
これは幾何学的には自然だが，コンパクト集合という位相的概念に依存している．
一般の集合 $X$ では位相を持たないことも多く，内測度を同じ形で定義できるとは限らない．

Carathéodoryの考え方は，内測度を使わず，外測度だけから可測性を判定することである．
この方法では「集合を $A$ の内側と外側に切っても外測度が余計に増えない」
という条件が，可測性の役割を果たす．

\section{可測空間と測度}

\begin{definition}[$\sigma$-加法族]\label{def:sigma-algebra}
集合 $X$ の部分集合族 $\calS\subset\calP(X)$ が
$X$ 上の $\sigma$-加法族であるとは，次を満たすことをいう．
\begin{itemize}
    \item $X\in\calS$．
    \item $A\in\calS$ ならば $X\setminus A\in\calS$．
    \item $\{A_n\}_{n=1}^{\infty}\subset\calS$ ならば
    $\bigcup_{n=1}^{\infty}A_n\in\calS$．
\end{itemize}
\end{definition}

\begin{remark}
$\sigma$-加法族は，補集合と可算和を取っても外へ出ない集合族である．
したがって可算共通部分や差集合についても閉じている．
\end{remark}

\begin{definition}[集合系が生成する $\sigma$-加法族]\label{def:generated-sigma-algebra}
任意の集合族 $\calC\subset\calP(X)$ に対し，
$\calC$ を含む最小の $\sigma$-加法族を $\sigma(\calC)$ と書き，
$\calC$ が生成する $\sigma$-加法族という．
具体的には，$\calC$ を含むすべての $\sigma$-加法族の共通部分である．
\end{definition}

\begin{definition}[可測空間]
$X$ 上の $\sigma$-加法族 $\calS$ を1つ固定した組 $(X,\calS)$ を可測空間という．
\end{definition}

\begin{definition}[測度]\label{def:measure}
可測空間 $(X,\calS)$ 上の写像 $\mu:\calS\to[0,\infty]$ が測度であるとは，
次を満たすことをいう．
\begin{itemize}
    \item $\mu(\emptyset)=0$．
    \item 互いに素な列 $\{A_n\}_{n=1}^{\infty}\subset\calS$ に対して
    $\mu(\bigsqcup_{n=1}^{\infty}A_n)=\sum_{n=1}^{\infty}\mu(A_n)$
    が成り立つ．
\end{itemize}
\end{definition}

\begin{definition}[測度空間]
可測空間 $(X,\calS)$ とその上の測度 $\mu$ の組 $(X,\calS,\mu)$ を測度空間という．
\end{definition}

\begin{definition}[$\sigma$-有限測度]\label{def:sigma-finite-measure}
測度 $\mu$ が $\sigma$-有限であるとは，
ある列 $\{X_n\}_{n=1}^{\infty}\subset\calS$ が存在して，
次を満たすことをいう．
\begin{itemize}
    \item $X=\bigcup_{n=1}^{\infty}X_n$．
    \item すべての $n$ について $\mu(X_n)<\infty$．
\end{itemize}
同じ言葉を有限加法族上の前測度に対して使うときも，
被覆集合 $X_n$ はその有限加法族の元とする．
\end{definition}

\begin{example}[数え上げ測度]\label{ex:counting-measure}
任意の集合 $X$ について $\calP(X)$ 上で
$\#A$ を $A$ の元の個数，ただし無限集合なら $\#A=\infty$，と定めると測度になる．
これを数え上げ測度（計数測度，個数測度）という．
$X$ が可算集合なら一点集合の可算和で覆えるので $\sigma$-有限であり，
$X$ が非可算集合なら $\sigma$-有限ではない．
\end{example}

\begin{example}[Dirac測度]\label{ex:dirac-measure}
可測空間 $(X,\calS)$ と点 $a\in X$ に対し，
$\delta_a(A)=1$（$a\in A$ のとき），$\delta_a(A)=0$（$a\notin A$ のとき）と定める．
これは $(X,\calS)$ 上の測度であり，点 $a$ に質量が集中したDirac測度という．
\end{example}

\section{外測度とCarathéodory可測性}

\begin{definition}[外測度]\label{def:outer-measure}
集合 $X$ 上の写像 $\mu^*:\calP(X)\to[0,\infty]$ が外測度であるとは，
次を満たすことをいう．
\begin{itemize}
    \item $\mu^*(\emptyset)=0$．
    \item $A\subset B$ ならば $\mu^*(A)\le\mu^*(B)$．
    \item 任意の集合列 $\{A_n\}_{n=1}^{\infty}\subset\calP(X)$ に対して
    $\mu^*(\bigcup_{n=1}^{\infty}A_n)\le\sum_{n=1}^{\infty}\mu^*(A_n)$．
\end{itemize}
\end{definition}

\begin{definition}[Carathéodory条件]\label{def:caratheodory-condition}
外測度 $\mu^*$ に対し，集合 $A\subset X$ がCarathéodory条件を満たすとは，
任意の $E\subset X$ に対して
\[
\mu^*(E)=\mu^*(E\cap A)+\mu^*(E\setminus A)
\]
が成り立つことをいう．
この条件を満たす集合を $\mu^*$-可測集合ともいう．
\end{definition}

\begin{theorem}[Carathéodoryの定理]\label{thm:caratheodory-theorem}
外測度 $\mu^*$ に関する $\mu^*$-可測集合全体を $\calM_{\mu^*}$ と書く．
このとき次が成り立つ．
\begin{itemize}
    \item $\calM_{\mu^*}$ は $X$ 上の $\sigma$-加法族である．
    \item $\mu:=\mu^*|_{\calM_{\mu^*}}$ は $(X,\calM_{\mu^*})$ 上の測度である．
    \item $\mu(N)=0$ である可測集合 $N$ の任意の部分集合も可測で，測度 $0$ をもつ．
\end{itemize}
\end{theorem}
\begin{proof}
外測度の劣加法性より，Carathéodory条件では
$\mu^*(E)\ge\mu^*(E\cap A)+\mu^*(E\setminus A)$ だけを示せばよい．
$\emptyset$ と $X$ が可測であること，および補集合で閉じていることは定義から直ちに従う．

まず有限和で閉じていることを示す．
$A,B$ を可測集合とし，任意の $E\subset X$ を取る．
$A$ の可測性より
\[
\mu^*(E)
=
\mu^*(E\cap A)+\mu^*(E\setminus A)
\]
である．
さらに $B$ の可測性を $E\setminus A$ に適用すると
\[
\mu^*(E\setminus A)
=
\mu^*((E\setminus A)\cap B)
+
\mu^*((E\setminus A)\setminus B)
\]
である．
ここで $(E\setminus A)\setminus B=E\setminus(A\cup B)$ だから
\[
\mu^*(E)
=
\mu^*(E\cap A)
+
\mu^*((E\setminus A)\cap B)
+
\mu^*(E\setminus(A\cup B)).
\]
一方，
\[
E\cap(A\cup B)
=
(E\cap A)\cup((E\setminus A)\cap B)
\]
なので，外測度の劣加法性より
\[
\mu^*(E\cap(A\cup B))
\le
\mu^*(E\cap A)+\mu^*((E\setminus A)\cap B)
\]
である．
したがって
\[
\mu^*(E)
\ge
\mu^*(E\cap(A\cup B))+\mu^*(E\setminus(A\cup B)).
\]
よって $A\cup B$ は可測である．
帰納法により有限和で閉じており，補集合で閉じているので差集合でも閉じている．

次に可算和で閉じていることを示す．
可測集合列 $\{A_n\}_{n=1}^{\infty}$ に対し，
\[
B_1:=A_1,\qquad
B_n:=A_n\setminus\bigcup_{k<n}A_k\quad(n\ge2)
\]
とおく．
有限和と差集合で閉じていることから各 $B_n$ は可測であり，
$B_n$ は互いに素で，
\[
\bigcup_{n=1}^{\infty}A_n=\bigsqcup_{n=1}^{\infty}B_n
\]
である．
$B:=\bigcup_{n=1}^{\infty}B_n$，$C_N:=\bigcup_{n=1}^{N}B_n$ とおく．
有限和で閉じているので $C_N$ は可測である．
また，互いに素な可測集合で順に切ると，任意の $E\subset X$ について
\[
\mu^*(E)
=
\sum_{n=1}^{N}\mu^*(E\cap B_n)
+
\mu^*(E\setminus C_N)
\]
が成り立つ．
実際，$N=1$ では $B_1$ の可測性そのものであり，
$N$ で成り立つときは $B_{N+1}$ の可測性を $E\setminus C_N$ に適用すればよい．
$E\setminus B\subset E\setminus C_N$ だから単調性より
\[
\mu^*(E)
\ge
\sum_{n=1}^{N}\mu^*(E\cap B_n)
+
\mu^*(E\setminus B)
\]
である．
$N\to\infty$ として
\[
\mu^*(E)
\ge
\sum_{n=1}^{\infty}\mu^*(E\cap B_n)
+
\mu^*(E\setminus B)
\]
を得る．
さらに外測度の可算劣加法性より
\[
\mu^*(E\cap B)
=
\mu^*\left(\bigcup_{n=1}^{\infty}(E\cap B_n)\right)
\le
\sum_{n=1}^{\infty}\mu^*(E\cap B_n)
\]
であるから，
\[
\mu^*(E)
\ge
\mu^*(E\cap B)+\mu^*(E\setminus B)
\]
となる．
よって $B$ は可測であり，可測集合の可算和も可測である．
よって $\calM_{\mu^*}$ は $\sigma$-加法族である．

互いに素な可測列 $A_n$ と $A=\bigsqcup_n A_n$ に対し，
上の有限段階の分解式を $E=A$，$B_n=A_n$ に適用すると
\[
\mu^*(A)
=
\sum_{n=1}^{N}\mu^*(A_n)
+
\mu^*\left(A\setminus\bigcup_{n=1}^{N}A_n\right)
\ge
\sum_{n=1}^{N}\mu^*(A_n)
\]
がすべての $N$ で成り立つ．
よって $\mu^*(A)\ge\sum_{n=1}^{\infty}\mu^*(A_n)$ である．
逆向きは外測度の可算劣加法性であるから，
$\mu(A)=\sum_n\mu(A_n)$ であり，$\mu$ は測度である．
最後に $N$ が可測で $\mu(N)=0$，$A\subset N$ とする．
任意の $E$ で $\mu^*(E\cap A)=0$ かつ $\mu^*(E\setminus A)\le\mu^*(E)$ なので，
$A$ もCarathéodory条件を満たし，測度は $0$ である．
\end{proof}

\begin{definition}[完備測度と完備化]\label{def:completion}
\begin{itemize}
    \item
    測度空間 $(X,\calS,\mu)$ が完備であるとは，
    $N\in\calS$，$\mu(N)=0$，$A\subset N$ ならば $A\in\calS$ であることをいう．
    \item
    測度（空間）の完備化とは，
    零集合の部分集合をすべて可測集合として加えた測度空間であり，
    たとえば
    $\overline{\calS}:=\{A\triangle Z\mid A\in\calS,\ Z\subset N,\ N\in\calS,\ \mu(N)=0\}$
    上に $\overline\mu(A\triangle Z):=\mu(A)$ と定めて得られる．
\end{itemize}
ここで $A\triangle Z:=(A\setminus Z)\cup(Z\setminus A)$ は対称差を表す．
\end{definition}

\section{半加法族と前測度}

\begin{definition}[半加法族]\label{def:semiring}
集合 $X$ の部分集合族 $\calE\subset\calP(X)$ が半加法族（半環）であるとは，
次を満たすことをいう．
\begin{itemize}
    \item $\emptyset\in\calE$．
    \item $A,B\in\calE$ ならば $A\cap B\in\calE$．
    \item $A,B\in\calE$ ならば $A\setminus B$ が有限個の互いに素な $\calE$ の元の合併
    として書ける．
\end{itemize}
ここで有限個は $0$ 個でもよい．
\end{definition}

\begin{remark}
$X\in\calE$ であっても，半加法族 $\calE$ は一般には有限加法族ではない．
半加法族の定義で要求するのは，差集合が有限個の $\calE$ の元の合併に分解できることまでであり，
その有限合併自体が再び $\calE$ に属するとは限らないからである．
\end{remark}

\begin{example}[左半開区間の半加法族]\label{ex:left-half-open-semiring}
$\RR^d$ では
\[
\calE_d
:=
\left\{
\prod_{i=1}^{d}(a_i,b_i]
\;\middle|\;
-\infty\le a_i<b_i\le\infty
\right\}
\cup\{\emptyset\}
\]
を用いる．
端点が $\pm\infty$ のときは自然な意味で解釈する．
共通部分は再び左半開区間であり，差集合は有限個の互いに素な左半開区間に分割できるので，
$\calE_d$ は半加法族である．
これまで基本集合と呼んできた族は $\calA(\calE_d)$ である．
\end{example}


\begin{example}[有限分割の半加法族]\label{ex:finite-partition-semiring}
$X$ の有限分割
\[
X=P_1\sqcup\cdots\sqcup P_m
\]
を固定する．
このとき
\[
\calE_{\Pi}:=\{\emptyset,P_1,\ldots,P_m\}
\]
は半加法族である．
実際，異なるブロックは互いに素であり，各ブロック同士の共通部分や差集合は
ブロックそのものか空集合になる．
ただし，$m\ge 2$ なら $P_1\cup P_2$ はふつう $\calE_{\Pi}$ に入らないので，
$\calE_{\Pi}$ は一般には有限加法族ではない．
生成される有限加法族 $\calA(\calE_{\Pi})$ は，分割のブロックの合併全体である．
\end{example}

\begin{example}[離散的な半加法族]\label{ex:discrete-semirings}
任意の集合 $X$ に対し，冪集合 $\calP(X)$ は半加法族である．
また，有限集合全体
\[
\calF(X):=\{F\subset X\mid F\text{ は有限集合}\}
\]
も半加法族である．
特に，$X$ が有限集合なら $\calF(X)=\calP(X)$ である．
これらは数え上げ測度，有限集合上の確率測度，Dirac測度を構成するときの自然な出発点になる．
\end{example}


\begin{definition}[有限加法族]\label{def:finite-additive-family}
集合 $X$ の部分集合族 $\calA\subset\calP(X)$ が有限加法族であるとは，
次を満たすことをいう．
\begin{itemize}
    \item $\emptyset\in\calA$．
    \item $A,B\in\calA$ ならば $A\cup B\in\calA$．
    \item $A,B\in\calA$ ならば $A\setminus B\in\calA$．
\end{itemize}
\end{definition}

\begin{definition}[半加法族が生成する有限加法族]\label{def:algebra-generated-by-semiring}
半加法族 $\calE$ に対して
\[
\calA(\calE)
:=
\left\{
\bigsqcup_{j=1}^{N}E_j
\;\middle|\;
N\in\NN,\ E_j\in\calE
\right\}
\]
と定める．
半加法族は $\emptyset$ を含むので，$\calA(\calE)$ も $\emptyset$ を含んでいる．
\end{definition}

\begin{theorem}\label{thm:algebra-generated-by-semiring}
$\calA(\calE)$ は $X$ 上の有限加法族である．
\end{theorem}
\begin{proof}
まず $C\in\calE$ と $F\in\calA(\calE)$ に対して，
$C\setminus F$ は有限個の互いに素な $\calE$ の元の合併になる．
これは $F$ の表示に現れる項数に関する帰納法と，半加法族の差集合条件から従う．
したがって $E=\bigsqcup_i E_i$ と $F\in\calA(\calE)$ なら
$E\setminus F=\bigsqcup_i(E_i\setminus F)$ も $\calA(\calE)$ に属する．
また $E\cup F=E\sqcup(F\setminus E)$ だから有限和でも閉じており，
$\emptyset$ は $\calE$ に属するので $\calA(\calE)$ にも属する．
\end{proof}

\begin{definition}[半加法族上の前測度]\label{def:premeasure-on-semiring}
半加法族 $\calE$ 上の写像 $\mu_0:\calE\to[0,\infty]$ が前測度であるとは，
次を満たすことをいう．
\begin{itemize}
    \item $\mu_0(\emptyset)=0$．
    \item 互いに素な列 $\{E_n\}_{n=1}^{\infty}\subset\calE$ が
    $E=\bigsqcup_{n=1}^{\infty}E_n\in\calE$ を満たすとき
    $\mu_0(E)=\sum_{n=1}^{\infty}\mu_0(E_n)$ となる．
\end{itemize}
半加法族上の前測度を体積と呼ぶこともある．
有限個の分割については，残りを空集合とすれば同じ等式が成り立つ．
\end{definition}

\begin{remark}
可算加法性は後の拡張定理において本質的な条件であるため，
有限加法性だけでなく，前測度の定義の時点で含めておくことが多い．
\end{remark}

\begin{remark}[用語の使い分け]
本章では，前測度を半加法族 $\calE$ から有限加法族 $\calA(\calE)$ へ移すことを
延長と呼ぶ．
また，前測度から外測度を誘導することを誘導と呼び，
得られた測度をより大きな $\sigma$-加法族上で考えることを拡張と呼ぶ．
\end{remark}

\begin{example}[離散的な前測度]\label{ex:discrete-premeasures}
\cref{ex:discrete-semirings}の $\calP(X)$ 上では，次はいずれも半加法族上の前測度である．
\begin{itemize}
    \item $\mu_0(A)=\#A$ と定めると，数え上げ測度の出発点になる．
    \item $X$ が有限集合で，$p_x\ge 0$，$\sum_{x\in X}p_x=1$ を満たす重みを固定し，
    $\mu_0(A)=\sum_{x\in A}p_x$ と定めると，有限集合上の確率測度の出発点になる．
    \item $a\in X$ を固定し，$\mu_0(A)=1$（$a\in A$ のとき），
    $\mu_0(A)=0$（$a\notin A$ のとき）と定めると，Dirac測度の出発点になる．
\end{itemize}
いずれも $\sigma(\calP(X))=\calP(X)$ なので，
拡張後の測度は同じ式で $\calP(X)$ 上に定まる．
\end{example}

\begin{definition}[有限加法族上の前測度]\label{def:premeasure-on-algebra}
有限加法族 $\calA$ 上の写像 $\mu_0:\calA\to[0,\infty]$ が前測度であるとは，
次を満たすことをいう．
\begin{itemize}
    \item $\mu_0(\emptyset)=0$．
    \item 互いに素な列 $\{A_n\}_{n=1}^{\infty}\subset\calA$ が
    $\bigsqcup_{n=1}^{\infty} A_n\in\calA$ を満たすとき
    $\mu_0(\bigsqcup_{n=1}^{\infty} A_n)=\sum_{n=1}^{\infty}\mu_0(A_n)$ となる．
\end{itemize}
\end{definition}

\begin{remark}
有限加法性はこの定義から従う．
実際，互いに素な有限列の後ろに空集合を並べれば，上の可算加法性を適用できる．
したがって，$A\subset B$ なら
$\mu_0(B)=\mu_0(A)+\mu_0(B\setminus A)$ となり，単調性も従う．
\end{remark}

\begin{theorem}[有限加法族への延長]\label{thm:premeasure-extension-to-algebra}
$\mu_0$ を半加法族 $\calE$ 上の前測度とする．
$A=\bigsqcup_{j=1}^{N}E_j\in\calA(\calE)$ に対して
$\overline\mu_0(A):=\sum_{j=1}^{N}\mu_0(E_j)$ と定める．
このとき次が成り立つ．
\begin{itemize}
    \item $\overline\mu_0$ は表示の取り方によらず定まる．
    \item $\overline\mu_0$ は $\calA(\calE)$ 上の前測度である．
\end{itemize}
\end{theorem}
\begin{proof}
$A=\bigsqcup_i E_i=\bigsqcup_j F_j$ と2通りに表す．
各 $E_i$ は互いに素な有限和 $E_i=\bigsqcup_j(E_i\cap F_j)$ に分割され，
各 $F_j$ も $F_j=\bigsqcup_i(E_i\cap F_j)$ に分割される．
半加法族上の前測度の有限加法性を使えば，
$\sum_i\mu_0(E_i)=\sum_{i,j}\mu_0(E_i\cap F_j)=\sum_j\mu_0(F_j)$ となる．
よって表示によらず定まる．
有限加法性は表示を並べればよい．
可算加法性は，$A=\bigsqcup_n A_n\in\calA(\calE)$ をさらに
$A=\bigsqcup_i E_i$ と書き，各 $E_i$ を $A_n\cap E_i$ で分割して
半加法族上の前測度の可算加法性を適用すれば従う．
\end{proof}

\section{誘導外測度と拡張定理}

\begin{definition}[誘導外測度]\label{def:outer-measure-generated-by-premeasure}
有限加法族 $\calA$ 上の前測度 $\mu_0$ から，任意の $B\subset X$ に対して
\[
\mu^*(B)
:=
\inf\left\{
\sum_{n=1}^{\infty}\mu_0(A_n)
\;\middle|\;
B\subset\bigcup_{n=1}^{\infty}A_n,\ A_n\in\calA\ (n=1,2,\ldots)
\right\}
\]
と定める．
つまり，下限は $B$ を覆う $\calA$ の可算列すべてにわたって取る．
そのような被覆が存在しないときは，この下限を $\infty$ と約束する．
\end{definition}

\begin{remark}[半加法族から直接作る流儀]\label{rem:outer-measure-from-semiring-directly}
半加法族 $\calE$ 上の前測度 $\mu_0$ から直接
\[
\mu^*_{\calE}(B)
:=
\inf\left\{
\sum_{n=1}^{\infty}\mu_0(E_n)
\;\middle|\;
B\subset\bigcup_{n=1}^{\infty}E_n,\ E_n\in\calE\ (n=1,2,\ldots)
\right\}
\]
と定める流儀もある．
一方，本章ではまず $\mu_0$ を有限加法族 $\calA(\calE)$ 上の前測度
$\overline\mu_0$ に延長し，その後で可算被覆による誘導外測度を作る．
両者は同じ外測度を与える．
実際，$\calE\subset\calA(\calE)$ なので $\calE$ による被覆は
$\calA(\calE)$ による被覆でもある．
逆に，$\calA(\calE)$ の各元は有限個の互いに素な $\calE$ の元の合併として書け，
そのとき $\overline\mu_0$ は各項の $\mu_0$ の和であるから，
$\calA(\calE)$ による可算被覆は同じ総和をもつ $\calE$ による可算被覆に細分できる．
したがって，下限は一致する．
\end{remark}

\begin{theorem}[誘導外測度]\label{thm:outer-measure-from-premeasure}
\cref{def:outer-measure-generated-by-premeasure}で定めた $\mu^*$ について，
次が成り立つ．
\begin{itemize}
    \item $\mu^*$ は $X$ 上の外測度である．
    \item 任意の $A\in\calA$ に対して $\mu^*(A)=\mu_0(A)$ である．
    \item 任意の $A\in\calA$ は $\mu^*$-可測である．
\end{itemize}
\end{theorem}
\begin{proof}
$\mu^*(\emptyset)=0$ は空集合による被覆から従い，単調性は被覆の範囲が広がるだけである．
可算劣加法性は，右辺が有限の場合に各 $E_n$ を $\eps 2^{-n}$ 以内で覆う
$\calA$ の列を取り，それらを全部並べることで従う．
次に $A\in\calA$ なら1項被覆で $\mu^*(A)\le\mu_0(A)$ である．
逆に $A\subset\bigcup_n A_n$ を $\calA$ の被覆とし，
$B_1=A\cap A_1$，$B_n=A\cap(A_n\setminus\bigcup_{k<n}A_k)$ とおくと，
$A=\bigsqcup_n B_n$，$B_n\in\calA$，$B_n\subset A_n$ である．
前測度性と単調性より $\mu_0(A)=\sum_n\mu_0(B_n)\le\sum_n\mu_0(A_n)$ なので，
下限を取って $\mu_0(A)\le\mu^*(A)$ を得る．
最後に $A\in\calA$ と任意の $E\subset X$ を取る．
$E\subset\bigcup_n A_n$ という任意の $\calA$-被覆を
$A_n\cap A$ と $A_n\setminus A$ に分けると，
$\mu^*(E\cap A)+\mu^*(E\setminus A)\le\sum_n\mu_0(A_n)$ である．
下限を取ればCarathéodory条件の逆向きが従い，劣加法性と合わせて $A$ は可測である．
\end{proof}

\begin{theorem}[Carathéodory--Hahnの拡張定理]\label{thm:caratheodory-hahn-extension}
$\calE$ を集合 $X$ 上の半加法族，
$\mu_0$ をその上の前測度とする．
$\overline\mu_0$ を $\calA(\calE)$ 上に延長した前測度，
$\mu^*$ を $\overline\mu_0$ から誘導した外測度，
$\calM_{\mu^*}$ を $\mu^*$-可測集合全体，
$\mu:=\mu^*|_{\calM_{\mu^*}}$ とする．
このとき次が成り立つ．
\begin{itemize}
    \item $\calE\subset\calS\subset\calM_{\mu^*}$ を満たす任意の $\sigma$-加法族
    $\calS$ に対し，$\mu_{\calS}:=\mu|_{\calS}$ は
    $(X,\calS)$ 上の測度であり，すべての $E\in\calE$ に対して
    $\mu_{\calS}(E)=\mu_0(E)$ を満たす．
    \item とくに $\sigma(\calE)=\sigma(\calA(\calE))\subset\calM_{\mu^*}$ であり，
    $\calS=\sigma(\calE)$ と取れば通常の生成 $\sigma$-加法族上の拡張が得られる．
    また，$\calS=\calM_{\mu^*}$ と取れば
    \cref{thm:caratheodory-theorem}で得られる測度空間
    $(X,\calM_{\mu^*},\mu)$ そのものである．
    \item $\overline\mu_0$ が $\calA(\calE)$ 上で $\sigma$-有限ならば，
    上の各 $\calS$ 上でこの拡張は一意的である．
\end{itemize}
\end{theorem}


\begin{proof}
\cref{thm:premeasure-extension-to-algebra}により
$\overline\mu_0$ は $\calA(\calE)$ 上の前測度である．
\cref{thm:outer-measure-from-premeasure}で外測度 $\mu^*$ を作ると，
$\calA(\calE)$ の元はすべて $\mu^*$-可測で，$\mu^*$ はその上で $\overline\mu_0$ と一致する．
\cref{thm:caratheodory-theorem}より
$\mu=\mu^*|_{\calM_{\mu^*}}$ は測度であり，
$\sigma(\calE)=\sigma(\calA(\calE))\subset\calM_{\mu^*}$ である．
したがって $\calE\subset\calS\subset\calM_{\mu^*}$ を満たす任意の $\calS$ について，
$\mu_{\calS}:=\mu|_{\calS}$ は $(X,\calS)$ 上の測度である．
また $E\in\calE$ なら
$\mu_{\calS}(E)=\mu^*(E)=\overline\mu_0(E)=\mu_0(E)$ であるから，これは求める拡張である．

一意性を示す．
$(X,\calS)$ 上の別の拡張を $\nu$ とする．
$\calA(\calE)\subset\calS$ であり，有限加法性により
$\nu$ と $\mu_{\calS}$ は $\calA(\calE)$ 上で一致する．
$\sigma$-有限性により $X_n\in\calA(\calE)$，$X_n\uparrow X$，
$\overline\mu_0(X_n)<\infty$ と取る．
任意の $F\in\calS$ に対し $F_n:=F\cap X_n$ とおく．
任意の $\calA(\calE)$-被覆 $F_n\subset\bigcup_k A_k$ について
$\nu(F_n)\le\sum_k\overline\mu_0(A_k)$ だから $\nu(F_n)\le\mu^*(F_n)$．
逆向きは，$X_n\setminus F$ の任意の $\calA(\calE)$-被覆を用いて
$\overline\mu_0(X_n)-\nu(F_n)\le\mu^*(X_n\setminus F)$ とし，
$X_n$ と $F$ の $\mu^*$-可測性から
$\mu^*(F_n)=\overline\mu_0(X_n)-\mu^*(X_n\setminus F)\le\nu(F_n)$ を得る．
よって $\nu(F_n)=\mu_{\calS}(F_n)$ であり，$F_n\uparrow F$ と測度の下からの連続性により
$\nu(F)=\mu_{\calS}(F)$ である．
\end{proof}

\begin{remark}[Hopfの拡張定理との関係]
用語には揺れがあるが，ここでは前測度から測度が存在する部分を
Carathéodory--Hahnの拡張定理と呼んだ．
$\sigma$-有限性の仮定のもとで一意性まで述べる形をHopfの拡張定理と呼ぶ文献もある．
要点は，半加法族 $\calE$ から有限加法族 $\calA(\calE)$ へ延長し，
そこから外測度を誘導し，Carathéodory可測集合へ制限するという流れである．
\end{remark}

\begin{example}[Lebesgue測度]\label{ex:lebesgue-measure-from-caratheodory-hahn}
\cref{ex:left-half-open-semiring}の $\calE_d$ 上で
$\mu_0(\prod_{i=1}^d(a_i,b_i])=\prod_{i=1}^d(b_i-a_i)$，
$\mu_0(\emptyset)=0$ と定める．
これは半加法族上の前測度である．
この前測度から誘導される外測度はLebesgue外測度の定義そのものであり，
\cref{cor:lebesgue-caratheodory-equivalence}により
Carathéodory可測集合全体はLebesgue可測集合族 $\calL_d$ と一致する．
したがって \cref{thm:caratheodory-hahn-extension} を
$\calS=\calL_d$ に適用すると，$\calL_d$ 上に得られる測度
$\mu^*|_{\calL_d}$ はLebesgue測度そのものである．
\end{example}

% \begin{remark}[Riemann積分と有限分割]\label{rem:riemann-partitions-and-semirings}
% 有界区間 $[a,b]$ の分割
% %\[
% $a=x_0<x_1<\cdots<x_n=b$
% %\]
% を固定する．
% Riemann積分では閉区間 $[x_{i-1},x_i]$ を考えるが，
% これは隣同士が端点で重なる．
% そこで
% \[
% I_1:=[x_0,x_1],
% \qquad
% I_i:=(x_{i-1},x_i]\quad (2\le i\le n)
% \]
% と互いに素なセルに直す．
% このとき
% \[
% \calE_{\Delta}:=\{\emptyset,I_1,\ldots,I_n\}
% \]
% は半加法族である．
% したがってRiemann和で用いる区間分割は，
% 半加法族から有限加法族へ進む構成の有限版と見ることができる．
% 分割をすべて動かすと，この型の半開小区間全体も半加法族になる．
%
% ただし，Riemann積分やDarboux積分は，
% この半加法族からCarathéodory外測度を作る積分論そのものではない．
% Riemann--Darbouxの議論では，分割の細分による有限個のセル上の和を考え，
% 上限や下限を有限分割全体にわたって取る．
% これは有限加法的なJordan内容に基づく積分論であり，
% Jordan外測度も有限被覆を使う量であって可算加法的な測度ではないため，
% 可算被覆から外測度を作って $\sigma$-加法的な測度へ拡張するCarathéodoryの手順とは別物である．
% Lebesgue積分論は，区間の長さを可算加法的な前測度として扱い，
% 可算被覆を許してCarathéodory外測度を作ることで，
% このRiemann--Darbouxの有限分割の議論をより広い測度論へ拡張している．
% \end{remark}

\begin{remark}[Jordan測度]\label{rem:jordan-caratheodory-branch}
有界区間（または有界矩形）$I$ に含まれる左半開区間全体に空集合を加えた半加法族を
$\calE(I)$ とする．
区間の長さ・体積は $\calE(I)$ 上の前測度を定める．
さらに，これを有限加法的に $\calA(\calE(I))$ へ延長するところまでは，
Lebesgue測度の構成とJordan測度の構成に共通する．

分岐点は，この有限加法族から任意の集合を外側近似する段階にある．
Carathéodory流では $\calA(\calE(I))$ の元による可算被覆を許して誘導外測度を作る．
一方，Jordan流では有限被覆だけを許す．
有限個の基本集合の合併は再び $\calA(\calE(I))$ に属するので，
Jordan外容量は
\[
\inf\{\overline\mu_0(G)\mid G\in\calA(\calE(I)),\ A\subset G\}
\]
という外側近似になる．
この有限被覆に留める点が，可算被覆から測度を作るCarathéodory流との違いである．
\end{remark}

\section{Borel集合族}

\begin{definition}[Borel集合族]\label{def:borel-sigma-algebra}
位相空間 $X$ に対し，$X$ の開集合全体を $\calO(X)$ と書く．
$\calB(X):=\sigma(\calO(X))$ を $X$ のBorel集合族といい，
$\calB(X)$ の元をBorel集合という．
\end{definition}

\begin{example}
$\RR^d$ では，開集合，閉集合，開集合の可算共通部分，閉集合の可算和はいずれもBorel集合である．
また，有理端点をもつ有界開区間
$\prod_{i=1}^d(a_i,b_i)$，$a_i,b_i\in\QQ$，の全体は可算基底をなすので，
$\calB(\RR^d)$ はこの可算な集合族から生成される．
\end{example}

\begin{theorem}[Borel集合はLebesgue可測]\label{thm:borel-sets-are-lebesgue-measurable}
すべてのBorel集合 $B\subset\RR^d$ はLebesgue可測である．
\end{theorem}
\begin{proof}
\cref{thm:open-closed-lebesgue-measurable}により開集合はLebesgue可測である．
また \cref{thm:lebesgue-measurable-sigma-algebra} により
Lebesgue可測集合全体 $\calL_d$ は $\sigma$-加法族である．
したがって開集合が生成する $\calB(\RR^d)$ は $\calL_d$ に含まれる．
\end{proof}

\begin{definition}[Borel測度]\label{def:borel-measure}
本章では，Lebesgue測度をBorel集合族に制限した測度
$\lambda|_{\calB(\RR^d)}$ を $\RR^d$ 上のBorel測度という．
\end{definition}

\begin{theorem}[Lebesgue測度の構造定理（Borel測度の完備化）]\label{thm:lebesgue-measure-completion-of-borel}
次が成り立つ．
\begin{itemize}
    \item 任意のLebesgue可測集合 $A\subset\RR^d$ は，
    あるBorel集合 $B$ とあるLebesgue零集合 $N$ を用いて $A=B\triangle N$ と書ける．
    \item Lebesgue測度は，Borel測度の完備化である．
\end{itemize}
\end{theorem}
\begin{proof}
まず $\lambda(A)<\infty$ の場合を示す．
\cref{thm:lebesgue-measure-outer-regularity}と有限加法性により，各 $n$ で
$A\subset O_n$ かつ $\lambda(O_n\setminus A)<1/n$ となる開集合 $O_n$ が取れる．
$B:=\bigcap_n O_n$ とおくと $B$ はBorel集合で，$A\subset B$，
かつ $B\setminus A\subset O_n\setminus A$ がすべての $n$ で成り立つので
$\lambda(B\setminus A)=0$ である．
一般の場合は $Q_m:=(-m,m]^d$ とし，有限測度集合 $A\cap Q_m$ に上の議論を適用して，
Borel集合 $B_m$ で $A\cap Q_m\subset B_m\subset Q_m$ かつ
$\lambda(B_m\setminus A)=0$ となるものを取る．
$B:=\bigcup_m B_m$ とすれば $B$ はBorel集合，$A\subset B$，
$B\setminus A$ は零集合である．
$N:=B\setminus A$ とおけば $A=B\triangle N$ であり，これは
\cref{def:completion}の意味でBorel測度を完備化している．
なお任意のLebesgue零集合は，外正則性によりBorel零集合に含まれる．
\end{proof}

\section{Lebesgue可測性の限界}

ここまでで，Borel集合はすべてLebesgue可測であり
（\cref{thm:borel-sets-are-lebesgue-measurable}），
Lebesgue測度はBorel測度の完備化として得られることを見た
（\cref{thm:lebesgue-measure-completion-of-borel}）．
つまりLebesgue測度は，位相から自然に作られるBorel集合だけでなく，
零集合の任意の部分集合まで測れる．

しかし，選択公理を用いると，それでもなおLebesgue可測でない集合が存在する．
その標準的な例がVitali集合である．

\begin{theorem}[Vitali]\label{thm:vitali-nonmeasurable}
$[0,1]$ にはLebesgue非可測集合が存在する．
\end{theorem}
\begin{proof}
$[0,1]$ 上に $x\sim y \Longleftrightarrow x-y\in\QQ$ で同値関係を定める．
選択公理を用いて各同値類から1点ずつ選んだ集合を $V\subset[0,1]$ とする．
$q\in[-1,1]\cap\QQ$ に対し $V+q:=\{v+q\mid v\in V\}$ とおく．
もし $q_1\ne q_2$ なら $(V+q_1)\cap(V+q_2)=\emptyset$ である．
実際，$v_1+q_1=v_2+q_2$ なら $v_1-v_2=q_2-q_1\in\QQ$ なので
$v_1,v_2$ は同じ同値類に属し，代表元の選び方から $v_1=v_2$，
したがって $q_1=q_2$ となる．

また $[0,1]\subset\bigcup_{q\in[-1,1]\cap\QQ}(V+q)\subset[-1,2]$ である．
ここで $V$ がLebesgue可測であると仮定する．
\cref{prop:lebesgue-measure-translation-invariant}により各 $V+q$ も可測で
$\lambda(V+q)=\lambda(V)$ である．
\cref{thm:lebesgue-measure-countable-additivity}より
\[
\lambda\left(\bigcup_{q\in[-1,1]\cap\QQ}(V+q)\right)
=
\sum_{q\in[-1,1]\cap\QQ}\lambda(V)
\]
となる．
もし $\lambda(V)=0$ なら右辺は $0$ だが，左辺の集合は $[0,1]$ を含むので少なくとも $1$ である．
もし $\lambda(V)>0$ なら右辺は $\infty$ だが，左辺の集合は $[-1,2]$ に含まれるので高々 $3$ である．
どちらも矛盾するから，$V$ はLebesgue可測ではない．
\end{proof}

\begin{remark}
Lebesgue測度はBorel集合より広い集合族まで測れるが，
\cref{thm:vitali-nonmeasurable}により $\calP(\RR^d)$ 全体を測れるわけではない．
Vitali集合は，可算加法性と平行移動不変性を保ったまま
すべての部分集合へLebesgue測度を拡張できないことを示している．
\end{remark}

\begin{theorem}[Solovay]\label{thm:solovay-model}
到達不能基数の存在を含むZFCが無矛盾であるならば，
ZFに従属選択公理を加えた公理系，すなわちZF+DC，のモデルで，
$\RR$ のすべての部分集合がLebesgue可測であるものが存在する．
\end{theorem}

\begin{remark}
Solovayの定理は，Lebesgue非可測集合の存在が
集合論の公理，とくに選択公理の強さに深く関わっていることを示している．
標準的な選択公理を仮定する通常のZFCではVitali集合が作れるが，
選択公理を弱めた世界では「すべての実数集合がLebesgue可測」と矛盾しない場合がある．
ここでは証明せず，非可測集合の存在は単なる測度論内部の現象ではなく，
集合論的な現象でもある，という位置づけだけを押さえる．
\end{remark}

\begin{theorem}[Banach--Tarskiのパラドックス]\label{thm:banach-tarski}
$\RR^3$ の球は有限個の互いに素な部分集合に分割でき，
それらを回転と平行移動だけで組み替えることで，
元の球と同じ大きさの球を2つ作ることができる．
\end{theorem}

\begin{remark}
Banach--Tarskiのパラドックスも選択公理を用いる結果である．
分割片は通常の意味で体積をもつ集合ではないため，
Lebesgue測度の有限加法性や平行移動不変性とは矛盾しない．
むしろこの結果は，$\RR^3$ のすべての部分集合に対して，
通常の体積を拡張するような有限加法的かつ運動不変な測度を定めることができないことを強く示している．
\end{remark}

\begin{remark}[任意加法性は強すぎる]
Lebesgue測度は可算加法的であるが，
任意個の互いに素な集合族に対してまで加法性を要求することはできない．
実際，すべての1点集合は測度 $0$ であるから，
もし任意加法性
\[
\lambda\left(\bigsqcup_{\alpha \in \Gamma} A_\alpha\right)
=
\sum_{\alpha \in \Gamma} \lambda(A_\alpha)
\]
が常に成り立つとする
（右辺は有限部分和の上限として解釈する）．
すると
\[
[0,1] = \bigsqcup_{x \in [0,1]} \{x\}
\]
より
\[
\lambda_1([0,1]) = \sum_{x \in [0,1]} 0 = 0
\]
となってしまう．
これは明らかに矛盾である．
したがって，面積を正しくモデル化するには
「可算加法性」がちょうどよい強さである．
\end{remark}

\section{まとめ}

本章では，Lebesgue流の測度論をCarathéodory流の一般論として捉え直した．
内測度を使って定義したLebesgue可測性は，
\cref{cor:lebesgue-caratheodory-equivalence}により
Lebesgue外測度に関するCarathéodory可測性と一致する．
この観点では，可測集合とは外測度を内側と外側に正しく分割する集合であり，
\cref{thm:caratheodory-theorem}により，その全体は自然に $\sigma$-加法族をなす．

また，半加法族上の前測度から誘導外測度を作り，
Carathéodory可測集合へ制限することで測度を得る手順を見た．
\cref{thm:caratheodory-hahn-extension}は，
区間の体積からLebesgue測度を作る構成が，
一般の測度構成の一例であることを説明している．
この枠組みでは，数え上げ測度やDirac測度も同じ「可測空間上の測度」として扱える．

Borel集合族は位相から生成される標準的な $\sigma$-加法族であり，
Lebesgue測度はBorel測度の完備化として得られる．
一方で，Vitali集合やBanach--Tarskiのパラドックスが示すように，
可算加法性，平行移動不変性，通常の体積との整合性を保ったまま
すべての部分集合を測ることはできない．
これで，Lebesgue積分を支える測度論の土台は一段落する．
